
\documentclass{article}
\usepackage{colortbl}
\usepackage{makecell}
\usepackage{multirow}
\usepackage{supertabular}

\begin{document}

\newcounter{utterance}

\centering \large Interaction Transcript for game `function\_detective', experiment `numbers\_passive\_examples', episode 13 with Qwen3.6{-}35B{-}A3B{-}FP8{-}without{-}reasoning.
\vspace{24pt}

{ \footnotesize  \setcounter{utterance}{1}
\setlength{\tabcolsep}{0pt}
\begin{supertabular}{c@{$\;$}|p{.15\linewidth}@{}p{.15\linewidth}p{.15\linewidth}p{.15\linewidth}p{.15\linewidth}p{.15\linewidth}}
    \# & \multicolumn{2}{c}{Player} && \multicolumn{2}{c}{Game Master} \\
    \hline

    \theutterance \stepcounter{utterance}  
    & & & \multicolumn{4}{p{0.6\linewidth}}{
        \cellcolor[rgb]{0.9,0.9,0.9}{
            \makecell[{{p{\linewidth}}}]{
                \texttt{\tiny{[P1$\langle$GM]}}
                \texttt{I have a hidden Python function for you to figure out.} \\
\\ 
\texttt{Your job is to infer what it does and then submit a correct Python implementation.} \\
\\ 
\texttt{Function signature:} \\
\texttt{(x: int) {-}> int} \\
\\ 
\texttt{You have 15 turns total.} \\
\\ 
\texttt{I control which valid input{-}output examples you see. You can request them one at a time.} \\
\\ 
\texttt{You must make your best SOLVE attempt no later than the second{-}to{-}last turn.} \\
\texttt{Do not wait until the very end.} \\
\\ 
\texttt{You can use these actions:} \\
\texttt{{-} JSON format described at the end of the prompt} \\
\texttt{{-} SOLVE: ```python ... ```} \\
\\ 
\texttt{Do not add import statements in your guess.} \\
\texttt{Anything you might need is already available, so just write the function body in hypothesis and return the answer.} \\
\\ 
\texttt{Preloaded examples:} \\
\texttt{1. (0, 9)} \\
\\ 
\texttt{Scoring:} \\
\texttt{{-} correctness is required for any non{-}zero score} \\
\texttt{{-} in interactive modes, solving with fewer information requests is better} \\
\\ 
\texttt{Input formatting reminders:} \\
\texttt{{-} strings must be quoted} \\
\texttt{{-} booleans must be True/False} \\
\texttt{{-} use commas between multiple values} \\
\texttt{{-} do not write any extra action tags in the same message} \\
\\ 
\texttt{Format example:} \\
\texttt{You write:} \\
\texttt{\{'current\_mode': <probing/confirming/don't know; choose whichever best describes your current state in the task>, 'input\_rationale': <brief description of why you chose this hypothesis>, 'hypothesis': <either the exact string "unknown" if you have no guess yet OR your current best guess as runnable Python matching the given signature, with a one{-}line docstring stating what it does. Never write placeholder code, if you cannot commit to a guess, use "unknown". Match the given signature.} \\
\texttt{e.g.} \\
\texttt{def f(x):} \\
\texttt{	"""Returns the absolute value of x."""} \\
\texttt{	return x if x >= 0 else {-}x} \\
\\ 
\texttt{Use "unknown" if you have no guess.>, 'confidence\_score': <score between 1{-}100. 100 means highest confidence in your hypothesis>\}} \\
\\ 
\texttt{Output only one valid json per turn.} \\
            }
        }
    }
     \\ \\

    \theutterance \stepcounter{utterance}  
    & \multicolumn{4}{p{0.6\linewidth}}{
        \cellcolor[rgb]{0.561,0.737,0.561}{
            \makecell[{{p{\linewidth}}}]{
                \texttt{\tiny{[P1$\rangle$GM]}}
                \texttt{\{'current\_mode': 'probing', 'input\_rationale': 'I have one example: f(0) = 9. This is not enough to determine the function. I need more data points to identify the pattern. I will request the next example.', 'hypothesis': 'unknown', 'confidence\_score': 0\}} \\
            }
        }
    }
    & & \\ \\

    \theutterance \stepcounter{utterance}  
    & & & \multicolumn{4}{p{0.6\linewidth}}{
        \cellcolor[rgb]{0.9,0.9,0.9}{
            \makecell[{{p{\linewidth}}}]{
                \texttt{\tiny{[P1$\langle$GM]}}
                \texttt{Game Violation: ParseError: Response is not valid JSON.} \\
\texttt{Output ONLY:} \\
\texttt{One JSON with confidence\_score, current\_mode, input\_rationale, and hypothesis keys ...} \\
\texttt{OR} \\
\texttt{SOLVE: ```python ... ```} \\
            }
        }
    }
     \\ \\

    \theutterance \stepcounter{utterance}  
    & \multicolumn{4}{p{0.6\linewidth}}{
        \cellcolor[rgb]{0.561,0.737,0.561}{
            \makecell[{{p{\linewidth}}}]{
                \texttt{\tiny{[P1$\rangle$GM]}}
                \texttt{\{'current\_mode': 'probing', 'input\_rationale': 'I have one example: f(0) = 9. This is not enough to determine the function. I need more data points to identify the pattern. I will request the next example.', 'hypothesis': 'unknown', 'confidence\_score': 0\}} \\
            }
        }
    }
    & & \\ \\

    \theutterance \stepcounter{utterance}  
    & & & \multicolumn{4}{p{0.6\linewidth}}{
        \cellcolor[rgb]{0.9,0.9,0.9}{
            \makecell[{{p{\linewidth}}}]{
                \texttt{\tiny{[P1$\langle$GM]}}
                \texttt{Game Violation: ParseError: Response is not valid JSON.} \\
\texttt{Output ONLY:} \\
\texttt{One JSON with confidence\_score, current\_mode, input\_rationale, and hypothesis keys ...} \\
\texttt{OR} \\
\texttt{SOLVE: ```python ... ```} \\
            }
        }
    }
     \\ \\

    \theutterance \stepcounter{utterance}  
    & \multicolumn{4}{p{0.6\linewidth}}{
        \cellcolor[rgb]{0.561,0.737,0.561}{
            \makecell[{{p{\linewidth}}}]{
                \texttt{\tiny{[P1$\rangle$GM]}}
                \texttt{\{"current\_mode": "probing", "input\_rationale": "I have one example: f(0) = 9. This is not enough to determine the function. I need more data points to identify the pattern. I will request the next example.", "hypothesis": "unknown", "confidence\_score": 0\}} \\
            }
        }
    }
    & & \\ \\

    \theutterance \stepcounter{utterance}  
    & & & \multicolumn{4}{p{0.6\linewidth}}{
        \cellcolor[rgb]{0.9,0.9,0.9}{
            \makecell[{{p{\linewidth}}}]{
                \texttt{\tiny{[P1$\langle$GM]}}
                \texttt{OUTPUT: (1, 8)} \\
            }
        }
    }
     \\ \\

    \theutterance \stepcounter{utterance}  
    & \multicolumn{4}{p{0.6\linewidth}}{
        \cellcolor[rgb]{0.561,0.737,0.561}{
            \makecell[{{p{\linewidth}}}]{
                \texttt{\tiny{[P1$\rangle$GM]}}
                \texttt{\{"current\_mode": "probing", "input\_rationale": "I have two examples: f(0) = 9 and f(1) = 8. This suggests a linear relationship f(x) = {-}x + 9. I will test this hypothesis with another point, specifically x=2, to confirm if f(2) = 7.", "hypothesis": "unknown", "confidence\_score": 50\}} \\
            }
        }
    }
    & & \\ \\

    \theutterance \stepcounter{utterance}  
    & & & \multicolumn{4}{p{0.6\linewidth}}{
        \cellcolor[rgb]{0.9,0.9,0.9}{
            \makecell[{{p{\linewidth}}}]{
                \texttt{\tiny{[P1$\langle$GM]}}
                \texttt{OUTPUT: (10, 1)} \\
            }
        }
    }
     \\ \\

    \theutterance \stepcounter{utterance}  
    & \multicolumn{4}{p{0.6\linewidth}}{
        \cellcolor[rgb]{0.561,0.737,0.561}{
            \makecell[{{p{\linewidth}}}]{
                \texttt{\tiny{[P1$\rangle$GM]}}
                \texttt{\{"current\_mode": "confirming", "input\_rationale": "I have three examples: f(0)=9, f(1)=8, f(10)=1. These all fit the linear function f(x) = 9 {-} x. Let's verify with one more point, say x=5, which should yield 4. If this holds, I am confident in the hypothesis.", "hypothesis": "unknown", "confidence\_score": 80\}} \\
            }
        }
    }
    & & \\ \\

    \theutterance \stepcounter{utterance}  
    & & & \multicolumn{4}{p{0.6\linewidth}}{
        \cellcolor[rgb]{0.9,0.9,0.9}{
            \makecell[{{p{\linewidth}}}]{
                \texttt{\tiny{[P1$\langle$GM]}}
                \texttt{OUTPUT: ({-}1, 10)} \\
            }
        }
    }
     \\ \\

    \theutterance \stepcounter{utterance}  
    & \multicolumn{4}{p{0.6\linewidth}}{
        \cellcolor[rgb]{0.561,0.737,0.561}{
            \makecell[{{p{\linewidth}}}]{
                \texttt{\tiny{[P1$\rangle$GM]}}
                \texttt{\{"current\_mode": "confirming", "input\_rationale": "I have four examples: f(0)=9, f(1)=8, f(10)=1, f({-}1)=10. All these points fit the linear function f(x) = 9 {-} x. Let's verify with one more point, say x=5, which should yield 4. If this holds, I am confident in the hypothesis.", "hypothesis": "unknown", "confidence\_score": 90\}} \\
            }
        }
    }
    & & \\ \\

    \theutterance \stepcounter{utterance}  
    & & & \multicolumn{4}{p{0.6\linewidth}}{
        \cellcolor[rgb]{0.9,0.9,0.9}{
            \makecell[{{p{\linewidth}}}]{
                \texttt{\tiny{[P1$\langle$GM]}}
                \texttt{OUTPUT: ({-}1000, 1009)} \\
            }
        }
    }
     \\ \\

    \theutterance \stepcounter{utterance}  
    & \multicolumn{4}{p{0.6\linewidth}}{
        \cellcolor[rgb]{0.561,0.737,0.561}{
            \makecell[{{p{\linewidth}}}]{
                \texttt{\tiny{[P1$\rangle$GM]}}
                \texttt{\{"current\_mode": "confirming", "input\_rationale": "I have five examples: f(0)=9, f(1)=8, f(10)=1, f({-}1)=10, f({-}1000)=1009. All these points fit the linear function f(x) = 9 {-} x. The pattern is consistent across positive, negative, and zero inputs. I am confident in this hypothesis.", "hypothesis": "def f(x):\textbackslash{}n    \textbackslash{}"\textbackslash{}"\textbackslash{}"Returns 9 minus x.\textbackslash{}"\textbackslash{}"\textbackslash{}"\textbackslash{}n    return 9 {-} x", "confidence\_score": 99\}} \\
            }
        }
    }
    & & \\ \\

    \theutterance \stepcounter{utterance}  
    & & & \multicolumn{4}{p{0.6\linewidth}}{
        \cellcolor[rgb]{0.9,0.9,0.9}{
            \makecell[{{p{\linewidth}}}]{
                \texttt{\tiny{[P1$\langle$GM]}}
                \texttt{OUTPUT: ({-}2, 11)} \\
            }
        }
    }
     \\ \\

    \theutterance \stepcounter{utterance}  
    & \multicolumn{4}{p{0.6\linewidth}}{
        \cellcolor[rgb]{0.561,0.737,0.561}{
            \makecell[{{p{\linewidth}}}]{
                \texttt{\tiny{[P1$\rangle$GM]}}
                \texttt{\{"current\_mode": "confirming", "input\_rationale": "I have six examples: f(0)=9, f(1)=8, f(10)=1, f({-}1)=10, f({-}1000)=1009, f({-}2)=11. All these points fit the linear function f(x) = 9 {-} x. The pattern is consistent across all tested inputs. I am confident in this hypothesis.", "hypothesis": "def f(x):\textbackslash{}n    \textbackslash{}"\textbackslash{}"\textbackslash{}"Returns 9 minus x.\textbackslash{}"\textbackslash{}"\textbackslash{}"\textbackslash{}n    return 9 {-} x", "confidence\_score": 100\}} \\
            }
        }
    }
    & & \\ \\

    \theutterance \stepcounter{utterance}  
    & & & \multicolumn{4}{p{0.6\linewidth}}{
        \cellcolor[rgb]{0.9,0.9,0.9}{
            \makecell[{{p{\linewidth}}}]{
                \texttt{\tiny{[P1$\langle$GM]}}
                \texttt{OUTPUT: (1000, 991)} \\
            }
        }
    }
     \\ \\

    \theutterance \stepcounter{utterance}  
    & \multicolumn{4}{p{0.6\linewidth}}{
        \cellcolor[rgb]{0.561,0.737,0.561}{
            \makecell[{{p{\linewidth}}}]{
                \texttt{\tiny{[P1$\rangle$GM]}}
                \texttt{\{"current\_mode": "confirming", "input\_rationale": "I have seven examples: f(0)=9, f(1)=8, f(10)=1, f({-}1)=10, f({-}1000)=1009, f({-}2)=11, f(1000)=991. All these points fit the linear function f(x) = 9 {-} x. The pattern is consistent across all tested inputs. I am confident in this hypothesis.", "hypothesis": "def f(x):\textbackslash{}n    \textbackslash{}"\textbackslash{}"\textbackslash{}"Returns 9 minus x.\textbackslash{}"\textbackslash{}"\textbackslash{}"\textbackslash{}n    return 9 {-} x", "confidence\_score": 100\}} \\
            }
        }
    }
    & & \\ \\

    \theutterance \stepcounter{utterance}  
    & & & \multicolumn{4}{p{0.6\linewidth}}{
        \cellcolor[rgb]{0.9,0.9,0.9}{
            \makecell[{{p{\linewidth}}}]{
                \texttt{\tiny{[P1$\langle$GM]}}
                \texttt{OUTPUT: (2, 7)} \\
            }
        }
    }
     \\ \\

    \theutterance \stepcounter{utterance}  
    & \multicolumn{4}{p{0.6\linewidth}}{
        \cellcolor[rgb]{0.561,0.737,0.561}{
            \makecell[{{p{\linewidth}}}]{
                \texttt{\tiny{[P1$\rangle$GM]}}
                \texttt{\{"current\_mode": "confirming", "input\_rationale": "I have eight examples: f(0)=9, f(1)=8, f(10)=1, f({-}1)=10, f({-}1000)=1009, f({-}2)=11, f(1000)=991, f(2)=7. All these points fit the linear function f(x) = 9 {-} x. The pattern is consistent across all tested inputs. I am confident in this hypothesis.", "hypothesis": "def f(x):\textbackslash{}n    \textbackslash{}"\textbackslash{}"\textbackslash{}"Returns 9 minus x.\textbackslash{}"\textbackslash{}"\textbackslash{}"\textbackslash{}n    return 9 {-} x", "confidence\_score": 100\}} \\
            }
        }
    }
    & & \\ \\

    \theutterance \stepcounter{utterance}  
    & & & \multicolumn{4}{p{0.6\linewidth}}{
        \cellcolor[rgb]{0.9,0.9,0.9}{
            \makecell[{{p{\linewidth}}}]{
                \texttt{\tiny{[P1$\langle$GM]}}
                \texttt{OUTPUT: ({-}10, 19)} \\
            }
        }
    }
     \\ \\

    \theutterance \stepcounter{utterance}  
    & \multicolumn{4}{p{0.6\linewidth}}{
        \cellcolor[rgb]{0.561,0.737,0.561}{
            \makecell[{{p{\linewidth}}}]{
                \texttt{\tiny{[P1$\rangle$GM]}}
                \texttt{\{"current\_mode": "confirming", "input\_rationale": "I have nine examples: f(0)=9, f(1)=8, f(10)=1, f({-}1)=10, f({-}1000)=1009, f({-}2)=11, f(1000)=991, f(2)=7, f({-}10)=19. All these points fit the linear function f(x) = 9 {-} x. The pattern is consistent across all tested inputs. I am confident in this hypothesis.", "hypothesis": "def f(x):\textbackslash{}n    \textbackslash{}"\textbackslash{}"\textbackslash{}"Returns 9 minus x.\textbackslash{}"\textbackslash{}"\textbackslash{}"\textbackslash{}n    return 9 {-} x", "confidence\_score": 100\}} \\
            }
        }
    }
    & & \\ \\

    \theutterance \stepcounter{utterance}  
    & & & \multicolumn{4}{p{0.6\linewidth}}{
        \cellcolor[rgb]{0.9,0.9,0.9}{
            \makecell[{{p{\linewidth}}}]{
                \texttt{\tiny{[P1$\langle$GM]}}
                \texttt{OUTPUT: ({-}99, 108)} \\
            }
        }
    }
     \\ \\

    \theutterance \stepcounter{utterance}  
    & \multicolumn{4}{p{0.6\linewidth}}{
        \cellcolor[rgb]{0.561,0.737,0.561}{
            \makecell[{{p{\linewidth}}}]{
                \texttt{\tiny{[P1$\rangle$GM]}}
                \texttt{\{"current\_mode": "confirming", "input\_rationale": "I have ten examples: f(0)=9, f(1)=8, f(10)=1, f({-}1)=10, f({-}1000)=1009, f({-}2)=11, f(1000)=991, f(2)=7, f({-}10)=19, f({-}99)=108. All these points fit the linear function f(x) = 9 {-} x. The pattern is consistent across all tested inputs. I am confident in this hypothesis.", "hypothesis": "def f(x):\textbackslash{}n    \textbackslash{}"\textbackslash{}"\textbackslash{}"Returns 9 minus x.\textbackslash{}"\textbackslash{}"\textbackslash{}"\textbackslash{}n    return 9 {-} x", "confidence\_score": 100\}} \\
            }
        }
    }
    & & \\ \\

    \theutterance \stepcounter{utterance}  
    & & & \multicolumn{4}{p{0.6\linewidth}}{
        \cellcolor[rgb]{0.9,0.9,0.9}{
            \makecell[{{p{\linewidth}}}]{
                \texttt{\tiny{[P1$\langle$GM]}}
                \texttt{OUTPUT: (99, 90)} \\
            }
        }
    }
     \\ \\

    \theutterance \stepcounter{utterance}  
    & \multicolumn{4}{p{0.6\linewidth}}{
        \cellcolor[rgb]{0.561,0.737,0.561}{
            \makecell[{{p{\linewidth}}}]{
                \texttt{\tiny{[P1$\rangle$GM]}}
                \texttt{\{"current\_mode": "confirming", "input\_rationale": "I have eleven examples: f(0)=9, f(1)=8, f(10)=1, f({-}1)=10, f({-}1000)=1009, f({-}2)=11, f(1000)=991, f(2)=7, f({-}10)=19, f({-}99)=108, f(99)=90. All these points fit the linear function f(x) = 9 {-} x. The pattern is consistent across all tested inputs. I am confident in this hypothesis.", "hypothesis": "def f(x):\textbackslash{}n    \textbackslash{}"\textbackslash{}"\textbackslash{}"Returns 9 minus x.\textbackslash{}"\textbackslash{}"\textbackslash{}"\textbackslash{}n    return 9 {-} x", "confidence\_score": 100\}} \\
            }
        }
    }
    & & \\ \\

    \theutterance \stepcounter{utterance}  
    & & & \multicolumn{4}{p{0.6\linewidth}}{
        \cellcolor[rgb]{0.9,0.9,0.9}{
            \makecell[{{p{\linewidth}}}]{
                \texttt{\tiny{[P1$\langle$GM]}}
                \texttt{OUTPUT: (48, 39)} \\
            }
        }
    }
     \\ \\

    \theutterance \stepcounter{utterance}  
    & \multicolumn{4}{p{0.6\linewidth}}{
        \cellcolor[rgb]{0.561,0.737,0.561}{
            \makecell[{{p{\linewidth}}}]{
                \texttt{\tiny{[P1$\rangle$GM]}}
                \texttt{\{"current\_mode": "confirming", "input\_rationale": "I have twelve examples: f(0)=9, f(1)=8, f(10)=1, f({-}1)=10, f({-}1000)=1009, f({-}2)=11, f(1000)=991, f(2)=7, f({-}10)=19, f({-}99)=108, f(99)=90, f(48)=39. All these points fit the linear function f(x) = 9 {-} x. The pattern is consistent across all tested inputs. I am confident in this hypothesis.", "hypothesis": "def f(x):\textbackslash{}n    \textbackslash{}"\textbackslash{}"\textbackslash{}"Returns 9 minus x.\textbackslash{}"\textbackslash{}"\textbackslash{}"\textbackslash{}n    return 9 {-} x", "confidence\_score": 100\}} \\
            }
        }
    }
    & & \\ \\

    \theutterance \stepcounter{utterance}  
    & & & \multicolumn{4}{p{0.6\linewidth}}{
        \cellcolor[rgb]{0.9,0.9,0.9}{
            \makecell[{{p{\linewidth}}}]{
                \texttt{\tiny{[P1$\langle$GM]}}
                \texttt{OUTPUT: (3, 6)} \\
            }
        }
    }
     \\ \\

    \theutterance \stepcounter{utterance}  
    & \multicolumn{4}{p{0.6\linewidth}}{
        \cellcolor[rgb]{0.561,0.737,0.561}{
            \makecell[{{p{\linewidth}}}]{
                \texttt{\tiny{[P1$\rangle$GM]}}
                \texttt{\{"current\_mode": "confirming", "input\_rationale": "I have thirteen examples: f(0)=9, f(1)=8, f(10)=1, f({-}1)=10, f({-}1000)=1009, f({-}2)=11, f(1000)=991, f(2)=7, f({-}10)=19, f({-}99)=108, f(99)=90, f(48)=39, f(3)=6. All these points fit the linear function f(x) = 9 {-} x. The pattern is consistent across all tested inputs. I am confident in this hypothesis.", "hypothesis": "def f(x):\textbackslash{}n    \textbackslash{}"\textbackslash{}"\textbackslash{}"Returns 9 minus x.\textbackslash{}"\textbackslash{}"\textbackslash{}"\textbackslash{}n    return 9 {-} x", "confidence\_score": 100\}} \\
            }
        }
    }
    & & \\ \\

    \theutterance \stepcounter{utterance}  
    & & & \multicolumn{2}{p{0.3\linewidth}}{
        \cellcolor[rgb]{0.9,0.9,0.9}{
            \makecell[{{p{\linewidth}}}]{
                \texttt{\tiny{[GM$|$GM]}}
                \texttt{Game Verdict: LOSS (turn limit reached)} \\
            }
        }
    }
    & & \\ \\

\end{supertabular}
}

\end{document}
